The October 2023 escalation generated a rapid surge of multilingual digital
discourse that researchers have begun to quantify computationally. Antonakaki
et al.~\cite{antonakaki2025} analyzed Telegram, Twitter/X, and Reddit from October
2023 to mid-2025 using LDA, BERTopic, and transformer-based emotion models,
identifying platform-specific narrative strategies and propaganda amplification
patterns. Santiago et al.~\cite{santiago2025} focused on Reddit, examining
80 days of conflict-related comments with ML-based sentiment and emotion
classifiers, and found that disgust dominated the emotional landscape. Closest
to our own design, Hofmann et al.~\cite{hofmann2026} hand-annotated 4{,}983
YouTube comments on the conflict for hate speech and for the same three stances
we use, trained classifiers reaching AUROC 0.83 to 0.90, and applied them to
comment sections from public and private news channels. Their corpus is roughly
two orders of magnitude smaller than ours and they do not compare platforms, but
their stance scheme and their finding of a predominantly neutral distribution
both anticipate ours. Notably, Antonakaki et al.\ already report that Reddit
hosts the more reflective and deliberative discussion of the platforms they
compare, which is one reason we treat that contrast as an established premise
rather than a finding of ours.

These studies establish that the conflict produced intense and polarized digital
discourse, but they leave three gaps this work addresses. First, neither includes
\textbf{YouTube}, so the media-centric reaction layer is absent from existing
cross-platform comparisons of this conflict. Second, neither assigns
\textbf{per-comment stance} at scale, which is what allows sentiment, toxicity and
interaction to be analysed \emph{conditional on} political position. Third, neither
examines \textbf{interaction structure} - reply networks, homophily, or
cross-platform stance transfer - so the question of whether these communities
actually talk to each other is untouched.

